% TOP_PROBLEM: {"id": 20000948, "title": "Simplex Ramsey numbers and a robust boundary-extension lemma", "category_id": 2, "category": {"id": 2, "name": "combinatorics", "display_name": "Combinatorics", "description": "Counting problems, graph theory, discrete structures.", "slug": "combinatorics", "order_index": 2, "created_at": "2026-07-31T15:26:25.671Z"}, "status": "partially_solved", "problem_number": "AIM-COMBINATORICS-0073", "set_id": 14}
% FIRST_POSTED: 2026-10-01
% ENTRY_KIND: statement_only
% UnsolvedMath ID 20000948; source catalogue code AIM-COMBINATORICS-0073
% !TeX program = lualatex
% Definition and sources only; this file is not an LLM solution attempt.
\documentclass[11pt,a4paper]{article}
\usepackage[margin=25mm]{geometry}
\usepackage{fontspec}
\setmainfont{lmroman10-regular.otf}[BoldFont=lmroman10-bold.otf,
 ItalicFont=lmroman10-italic.otf,BoldItalicFont=lmroman10-bolditalic.otf]
\usepackage{amsmath,amssymb,mathtools}
\usepackage{unicode-math}
\setmathfont{latinmodern-math.otf}
\AtBeginDocument{\let\setminus\smallsetminus}
\usepackage{xurl,hyperref}
\hypersetup{colorlinks=true,urlcolor=blue}
\setcounter{secnumdepth}{0}
\setlength{\parindent}{0pt}
\setlength{\parskip}{5pt}
\setlength{\emergencystretch}{3em}
\begin{document}
\section*{Simplex Ramsey numbers and a robust boundary-extension lemma}\label{20000948}
{\small UnsolvedMath ID 20000948 · AIM-COMBINATORICS-0073 · Combinatorics · AIM Workshop Problem Lists}
\subsection{Definitions and mathematical statement}
For $k\ge2$, a $k$-uniform hypergraph is a family of $k$-element
subsets of a vertex set. The complete such hypergraph on $k+1$
vertices, $K_{k+1}^{(k)}$, consists of all $k+1$ subsets obtained by
omitting one vertex. It is the hypergraph simplex.
Let $R_k$ be the least integer $N$ such that, whenever each member of
$\binom{[N]}k$ is colored red or blue, there is a set
$S\subseteq[N]$ of size $k+1$ for which all members of
$\binom Sk$ have the same color. Thus
\[
                        R_k=r_k(k+1,k+1).
\]
The question is to determine, or sharply estimate, the growth of $R_k$
as $k$ tends to infinity. Both the uniformity and the required clique
size grow; the clique has exactly one more vertex than the uniformity.
The number of colors remains two throughout.

An upper bound $R_k\le U(k)$ must force a monochromatic simplex in
every coloring on $U(k)$ vertices. A lower bound $R_k>L(k)$ requires
a coloring on $L(k)$ vertices with no such simplex, meaning every
$(k+1)$-vertex set has at least one edge of each color.
The desired estimates should clarify the growth scale in $k$, beyond
simply applying fixed-uniformity Ramsey finiteness. The vertices carry
no order restriction on copies, even though $[N]$ supplies labels.
\subsection{Short English statement}
How rapidly does the two-color Ramsey number of a hypergraph simplex grow with its uniformity?
\subsection{Sources}
\begin{itemize}
\item[S1] ULAM.AI and the UnsolvedMath Contributors, supplied problems.json, record 20000948 (AIM-COMBINATORICS-0073), AIM Workshop Problem Lists. Catalogue identity and source transcription; CC BY 4.0.
\url{https://huggingface.co/datasets/ulamai/UnsolvedMath}.
\item[S2] American Institute of Mathematics, Graph Ramsey theory, item 26. Original workshop question underlying AIM-COMBINATORICS-0073.
\url{https://www.overleaf.com/read/mnvcscjjysvg}.
\end{itemize}
\end{document}
